
\chapter{Metodología de Desarrollo, Validación y Despliegue}
\label{app:workflow-dev}

Una parte relevante del valor metodológico del proyecto es que no se ha construido como un prototipo monolítico difícil de auditar, sino como un repositorio vivo gestionado mediante prácticas habituales de ingeniería de software. El desarrollo se ha realizado sobre Git y GitHub, con versionado explícito del código, la configuración, la documentación, los artefactos de infraestructura y los componentes desplegables del sistema. Esta organización permite reconstruir la evolución del proyecto y asociar cada resultado técnico a una versión concreta del repositorio.

El flujo de trabajo se apoya en ramas de desarrollo, integración mediante \emph{pull requests} y automatización con GitHub Actions. De este modo, los cambios no se incorporan únicamente por modificación directa del código, sino a través de un proceso controlado donde cada entrega puede ser validada antes de integrarse en la rama principal. Esta práctica resulta especialmente importante en un proyecto que combina extracción y tratamiento de datos financieros, entrenamiento de modelos, explicabilidad, dashboard, API e infraestructura cloud, ya que una modificación local en cualquiera de estos bloques puede afectar silenciosamente al comportamiento global del sistema.

Para reforzar este control, el repositorio incorpora un workflow específico de validación sobre \emph{pull requests}. Dicho workflow instala el entorno Python del proyecto, ejecuta comprobaciones estáticas con \texttt{flake8} y lanza la batería de tests mediante \texttt{pytest}. Esto permite detectar errores de estilo, problemas de integración y regresiones funcionales antes de aceptar cambios en la rama principal. Además, el proyecto mantiene una cobertura de tests superior al 80\%, lo que aporta una garantía adicional sobre la estabilidad de los módulos principales y reduce el riesgo de introducir modificaciones no controladas en piezas críticas del sistema.

La cobertura de tests no debe interpretarse como una garantía absoluta de corrección, pero sí como una evidencia de madurez metodológica. En este trabajo, las pruebas automatizadas cubren componentes relevantes del repositorio y se complementan con validaciones manuales sobre los artefactos generados, los flujos de datos, el dashboard y la API. Esta combinación aproxima el desarrollo del proyecto al ciclo de vida de un proyecto real: implementar una mejora, verificar su comportamiento, integrarla de forma controlada y conservar una traza reproducible de los cambios realizados.

Junto al workflow de validación, el proyecto incluye automatización específica para la documentación. El workflow de despliegue de documentación construye el sitio generado con MkDocs y lo publica en GitHub Pages a partir de la configuración definida en \texttt{../doc/mkdocs.yml}. Esto permite que la documentación técnica evolucione junto con el código y que los cambios incorporados al repositorio puedan reflejarse de manera sistemática en una versión navegable del proyecto. La documentación deja así de ser un artefacto estático y separado, y pasa a formar parte del propio proceso de integración.

El repositorio también automatiza el despliegue de infraestructura cloud mediante Terraform. El workflow de infraestructura carga la configuración del proyecto desde \texttt{resources/cloud\_config.json}, prepara las variables necesarias para Terraform, autentica contra Google Cloud, inicializa el backend remoto de estado en Google Cloud Storage y ejecuta las fases de \texttt{plan} y \texttt{apply}. Este proceso permite definir la infraestructura como código, mantener trazabilidad sobre los recursos cloud utilizados y reducir la dependencia de configuraciones manuales no documentadas. Además, el uso de un backend remoto para el estado de Terraform facilita la consistencia entre ejecuciones y evita que el estado de la infraestructura dependa de una máquina local concreta.

Una vez preparada la infraestructura base, el sistema incorpora un workflow reutilizable para construir y desplegar los componentes ejecutables del proyecto: el dashboard y la API. Este workflow construye las imágenes Docker correspondientes a partir de los \texttt{Dockerfile} definidos para cada servicio, las publica en Artifact Registry y aplica la configuración necesaria para su despliegue en Cloud Run mediante Terraform. De esta forma, el paso desde código fuente hasta servicio desplegable queda formalizado en una cadena reproducible: construir imagen, publicar artefacto, planificar infraestructura y aplicar cambios.

Esta separación entre workflows tiene una finalidad metodológica clara. El workflow de \emph{pull request} protege la calidad del código antes de la integración; el workflow de documentación mantiene actualizada la documentación técnica; el workflow de infraestructura gestiona los recursos cloud de base; y el workflow de despliegue de dashboard y API materializa los servicios finales. Cada automatización cubre una responsabilidad distinta, evitando mezclar validación, documentación, provisión de infraestructura y despliegue aplicativo en un único proceso opaco.

En este contexto, operaciones como \texttt{commit}, \texttt{pull}, \texttt{push}, \emph{pull request}, \texttt{plan} o \texttt{apply} no son únicamente acciones operativas. Funcionan como mecanismos de trazabilidad metodológica. Un resultado experimental o una funcionalidad del sistema no queda defendida solo por su salida final, sino por la versión concreta del código, la configuración, los tests y la infraestructura que la acompañan. Esta trazabilidad es especialmente relevante en un proyecto financiero, donde pequeñas variaciones en el pipeline de datos, en la configuración de entrenamiento, en los modelos o en la capa de explicación pueden modificar los resultados obtenidos.

Por tanto, el uso de Git, GitHub Actions, tests automatizados, control estático de código, documentación desplegable, Docker, Terraform y Google Cloud no constituye un elemento accesorio de implementación. Forma parte de la metodología del trabajo: asegura reproducibilidad, facilita el control de cambios, reduce regresiones, documenta la evolución del sistema y permite justificar el resultado final no solo por lo que predice o visualiza, sino por cómo se ha construido, validado, integrado y desplegado.


